% !TeX program = xelatex
\documentclass[UTF8,fontset=none,zihao=-4,a4paper]{ctexart}
\usepackage[left=25mm,right=25mm,top=24mm,bottom=24mm,headheight=15pt]{geometry}
\usepackage{fontspec}
\setCJKmainfont{SimSun}[BoldFont=Microsoft YaHei,ItalicFont=KaiTi]
\setCJKsansfont{Microsoft YaHei}
\setCJKmonofont{Microsoft YaHei}
\setmainfont{Latin Modern Roman}
\usepackage{amsmath,amssymb,amsthm,mathtools,bm,booktabs,tabularx,longtable,array,enumitem,xcolor,hyperref,bookmark}
\definecolor{navy}{HTML}{173B55}
\hypersetup{unicode,colorlinks=true,linkcolor=navy,urlcolor=navy,citecolor=navy,pdftitle={PCP 的等价计算、最优性核验与重复尺度},pdfauthor={Anonymous research reading draft}}
\ctexset{section={format=\large\sffamily\bfseries\color{navy}},subsection={format=\normalsize\sffamily\bfseries}}
\newtheoremstyle{plainzh}{7pt}{7pt}{\normalfont}{}{\bfseries}{。}{.5em}{}
\theoremstyle{plainzh}
\newtheorem{proposition}{命题}[section]
\newtheorem{lemma}[proposition]{引理}
\newtheorem{example}[proposition]{例}
\newtheorem{remark}[proposition]{说明}
\renewcommand{\proofname}{证明}
\newcommand{\R}{\mathbb R}
\newcommand{\nuc}[1]{\lVert #1\rVert_*}
\newcommand{\op}[1]{\lVert #1\rVert_{\mathrm{op}}}
\newcommand{\fro}[1]{\lVert #1\rVert_F}
\newcommand{\ellone}[1]{\lVert #1\rVert_1}
\newcommand{\maxnorm}[1]{\lVert #1\rVert_\infty}
\newcommand{\ip}[2]{\langle #1,#2\rangle}
\newcommand{\prox}{\operatorname{prox}}
\newcommand{\SVT}{\mathcal D}
\newcommand{\soft}{\mathcal S}
\newcommand{\diag}{\operatorname{diag}}
\newcommand{\sgn}{\operatorname{sign}}
\newcommand{\rep}{\mathcal R}
\numberwithin{equation}{section}
\linespread{1.12}
\setlength{\emergencystretch}{2em}
\setlist{itemsep=3pt,topsep=4pt,leftmargin=2em}
\allowdisplaybreaks[2]
\title{\vspace{-1.4em}\bfseries PCP 的等价计算、最优性核验与重复尺度\\[.4em]\large 一个课程项目的数学与实验证据阅读稿}
\author{匿名整理工作稿}
\date{版本 1.0\quad 2026 年 11 月 25 日}
\begin{document}
\maketitle
\noindent\textbf{稿件状态。}本稿由课程小组项目引出，经 AI 辅助核读、推导和独立合成核验，尚待原作者复核。本文不主张提出新的 PCP、ADMM、SVT 或 BB 方法，也不作文献优先权判断。原始真实视频、原小组代码及组员身份不随本稿提供；随附的独立合成核验程序自行生成该程序使用的测试矩阵。
\begin{abstract}
鲁棒主成分分离的计算中，模型是否合适、近端是否算对以及迭代是否停止在可接受的点上，是三个不同的问题。本文围绕近端实现、停止规则和数据尺度之间的关系，给出一组可独立阅读的等价性与核验结果。首先推导 Gram 矩阵实现奇异值阈值算子的精确算术公式，并用一个三维例子说明部分奇异三元组的残差检查不能单独排除遗漏的活跃奇异值。随后从子问题最优性条件推导过松弛 ADMM 的正确 KKT 缺陷，配合全近端固定点检查与原始—对偶夹逼。对于严格重复的列块，证明最优值的缩放关系，并区分模型尺度变化和迭代轨道的参数变换。最后给出 Huber 消元及其 BB 割线步长的精确条件，结合原项目的全部实验组，划清算法等价、历史性能、模型误差与新增合成证据的边界。
\end{abstract}
\noindent\textbf{关键词：}主成分追踪；奇异值阈值；过松弛 ADMM；重复数据；Huber 损失；实验可核验性
\tableofcontents

\section{问题、材料与本文要回答的问题}
设 $D\in\R^{m\times n}$ 是观测矩阵，$\lambda>0$。主成分追踪（PCP）是凸优化问题
\begin{equation}\label{eq:pcp}
 v_\lambda(D)=\min_{L+S=D}\{\nuc L+\lambda\ellone S\}.
\end{equation}
这里 $\nuc L$ 是奇异值之和，$\ellone S$ 是全部元素绝对值之和，$\maxnorm S$ 是最大元素绝对值；$\op L$ 则是最大奇异值。内积统一为 $\ip AB=\operatorname{tr}(A^TB)$。对闭真凸函数 $h$，本文采用近端记号 $\prox_h(A)=\arg\min_X\{h(X)+\fro{X-A}^2/2\}$；子微分 $\partial h(X)$ 由所有满足 $h(W)\geq h(X)+\ip G{W-X}$ 的矩阵 $G$ 构成。将 $S$ 消去后，目标连续且随 $\fro L\to\infty$ 趋于无穷，因此最小值存在。存在最优解与该最优解恢复真实分量是两个不同结论：后一结论需要关于真实分量的结构条件，不能由凸性单独推出。PCP 的经典恢复理论见文献\cite{pcp}。

本次核读的主体是原项目最终报告（49 页）、课堂讲演（38 页）、复现说明以及算法、实验与复核程序。页码均指 PDF 的物理页码。各类算法和实验的对应关系列在附录 A。历史表格的核对、保存端点的复算及本次合成实验采用不同证据口径，下文分别说明。

需要区分两个资料范围。课程交付说明明确未在小包内打包超过 1 GB 的中间矩阵，课程目录与对应 ZIP 中均无完整端点的 \texttt{.mat} 文件；扩大私有资料范围后，另发现独立的大矩阵包及补遗矩阵。本稿的历史表核读先以小包为依据，随后已从独立私人包完成 8 组保存端点的重新计算，结果另列于第 3.2 节。它重算残差、近端和目标，但没有从初值重新运行全部优化过程。仅凭 CSV、图像或部分压缩数组仍不能重建完整的 $L,S,Y$。

已有独立专题处理了原报告第 2.4 节和附录 C 的常数背景、明暗矩形扰动及其对偶证书。本稿不重复这些条件的计算，而将余下内容组织成如下问题：
\begin{enumerate}
\item 用 Gram 特征分解替代 SVD，究竟等价到哪一层？浮点保护能证明什么？
\item 改变松弛参数后，何种残差仍然对应子问题最优性？怎样另行检查最终解？
\item 增加重复列会不会只是增加计算量？权重与罚参数应怎样解释？
\item 加入密集噪声后，Huber 模型及 BB 步长的表现能由什么机制解释？
\end{enumerate}
以下数学命题均给出证明。数值证据分为三类：原项目的历史实验记录（包括原合成实验）、已保存端点的新复核（包括真实视频及千帧合成序列），以及本次独立重新生成的合成核验。三者分别统计，不能相互替代。

\section{Gram 公式、完整近端与部分谱的边界}
\subsection{从核范数近端到 Gram 表达式}
对 $t\geq0$，记
\begin{equation}\label{eq:svt}
\SVT_t(A)=\prox_{t\|\cdot\|_*}(A)
=U\diag((\sigma_i-t)_+)V^T,
\end{equation}
其中 $A=U\diag(\sigma_i)V^T$ 是薄 SVD，$x_+=\max\{x,0\}$。当 $t=0$ 时约定 $\SVT_0(A)=A$。公式 \eqref{eq:svt} 是经典奇异值阈值公式\cite{svt}：由奇异值的迹不等式，固定候选奇异值时，二次距离在奇异向量与 $A$ 对齐时最小，余下的非负标量问题为
$\min_{x\geq0}\{tx+(x-\sigma_i)^2/2\}$，其唯一解是 $(\sigma_i-t)_+$。

\begin{proposition}[Gram 表达式；精确算术]\label{prop:gram}
若 $m\geq n$，令 $A^TA=V\diag(\sigma_i^2)V^T$，并令 $I=\{i:\sigma_i>t\}$。对 $t>0$，有
\begin{equation}\label{eq:gram}
 \SVT_t(A)=A V_I\diag\left(1-\frac{t}{\sigma_i}\right)_{i\in I}V_I^T.
\end{equation}
若 $m<n$，由 $AA^T$ 的特征向量 $U_I$ 得到
\begin{equation}
 \SVT_t(A)=U_I\diag\left(1-\frac{t}{\sigma_i}\right)_{i\in I}U_I^TA.
\end{equation}
空指标集对应零矩阵；重复奇异值、秩亏和阈值处的等号不破坏公式。
\end{proposition}
\begin{proof}
先看 $m\geq n$。对 $i\in I$，$\sigma_i>0$，定义 $u_i=Av_i/\sigma_i$。因为
$u_i^Tu_j=v_i^TA^TAv_j/(\sigma_i\sigma_j)$，这些向量正交归一，且 $Av_i=\sigma_i u_i$。于是 \eqref{eq:gram} 的右端是
$\sum_{i\in I}(\sigma_i-t)u_iv_i^T$，正好等于 \eqref{eq:svt}。在重复特征值的特征子空间内，乘子 $1-t/\sigma_i$ 相同，换取任何正交基均给出同一矩阵。$\sigma_i=t$ 时缩减后的系数为零，放入或移出求和都不改变结果。宽矩阵的论证相同，只需交换左右奇异向量。
\end{proof}

这个公式解释了原实现为什么没有改变 PCP 的数学子问题。它没有给出新的渐近复杂度：在 $m\geq n$ 时，构造 Gram 矩阵、特征分解及秩为 $r$ 的重构分别约为 $O(mn^2)$、$O(n^3)$ 和 $O(mnr)$；薄 SVD 也有同阶复杂度。能否更快取决于矩阵尺寸、线性代数库、内存访问和重构代价，需要实测。

\subsection{为什么浮点 Gram 仍需保守表述}
在非零奇异谱上，Gram 矩阵的条件数是 $A$ 的条件数的平方。小奇异值对应的平方因此可能落入舍入误差尺度；形成 $A^TA$ 还可能先于 $A$ 本身发生上溢或下溢。近阈值时，“保留了几项”也可能比近端矩阵本身更加敏感：一个仅比 $t$ 大 $\varepsilon$ 的奇异值只贡献大小为 $\varepsilon$ 的项。

有一个可严格陈述的条件式判据。若已经有可验证的误差界
\begin{equation}\label{eq:weyl}
 \op{\widehat G-A^TA}\leq\delta,
\end{equation}
并且 $\widehat G$ 对称，则按序排列的特征值满足
$|\widehat\mu_i-\sigma_i^2|\leq\delta$。因此，$\widehat\mu_i-\delta>t^2$ 足以保证活跃，$\widehat\mu_i+\delta<t^2$ 足以保证不活跃；若区间碰到阈值，单凭该界无法决定。证明是将 \eqref{eq:weyl} 代入对称矩阵的 Rayleigh 商，再用极小极大表征比较对应特征值。若计算的特征分解本身还带误差，该误差也必须计入 $\delta$。

原 \texttt{svt\_gram.m} 第 9 行已将保护量称为 heuristic；其第 67--88 行使用维度、机器精度、矩阵范数及特征残差构造保护。触及阈值、严重负特征值、非有限数值等情况会退回完整 SVD。这个设计是有根据的工程防护，但现有代码没有实现带向外舍入的 \eqref{eq:weyl} 证书。因此应写成“精确算术等价的 Gram 实现，附启发式浮点保护和完整 SVD 回退”，不能写成“所有输入上经过严格误差认证的精确 SVT”。完整 SVD 也使用浮点数，回退并不会把数值计算自动变成形式证明。

\subsection{部分奇异三元组为什么不够}
部分 SVD 的动机是只计算大于阈值的奇异值。若已经证明所有遗漏奇异值都不大于 $t$，它与完整 SVT 等价。难点正是这个“已经证明”。
\begin{example}[零残差仍可能漏掉活跃谱]\label{ex:partial}
取 $A=\diag(5,2,0.1)$、$t=1$。假设返回的两个奇异三元组为 $(e_2,2,e_2)$ 和 $(e_3,0.1,e_3)$。它们的左右残差都严格为零，向量正交，且最小返回奇异值 $0.1<t$；但据此重构只能得到 $\diag(0,1,0)$，而完整近端是 $\diag(4,1,0)$。二者 Frobenius 距离为 $4$。
\end{example}
这个例子没有声称原程序请求“最大奇异值”时实际发生了上述故障。它说明的是：局部三元组残差、正交性及最小返回值的检查，在逻辑上不能单独认证“返回的是全部应保留的大奇异值”。还需依赖求解器对目标谱的保证，或另有遗漏谱的上界。

原 \texttt{svt\_adaptive.m} 请求最大奇异值，检查求解器标志、三元组残差和正交性，按需扩展子空间，并在计算规模较大或检查失败时退回完整 SVD；源码也明确提示这些检查不是严格的遗漏谱证书。恰当的验证结构因此有两层：迭代内尽量减少近端成本；迭代后，用完整近端和可行的对偶变量重新检查最终结果。后者不能恢复前者每一步的精确性，却可以独立回答最终点是否满足所设精度。

\section{过松弛的 KKT 缺陷与端点核验}
\subsection{从两个子问题直接推导残差}
考虑更一般的 $\min\{f(L)+g(Z):L+Z=D\}$，其中 $f,g$ 为闭真凸函数。罚参数 $\rho>0$，乘子采用拉格朗日项 $\ip Y{L+Z-D}$ 的符号。固定一步，记旧变量为 $(L,Z,Y)$，令 $u=Y/\rho$，执行
\begin{align}
 L^+&=\prox_{f/\rho}(D-Z-u),\\
 \widehat L&=\alpha L^++(1-\alpha)(D-Z),\\
 Z^+&=\prox_{g/\rho}(D-\widehat L-u),\\
 Y^+&=Y+\rho(\widehat L+Z^+-D).
\end{align}
标准过松弛分析通常取 $0<\alpha<2$；下面的代数恒等式本身不依赖这个区间。设
\begin{equation}
 a=L^++Z-D,\qquad \Delta Z=Z^+-Z,\qquad r=L^++Z^+-D=a+\Delta Z.
\end{equation}
注意乘子更新使用的是 $\alpha a+\Delta Z$，而最终原始可行性使用的是 $r$。
\begin{proposition}[过松弛步的子梯度缺陷]\label{prop:relax}
若两个近端子问题均精确求解，则
\begin{equation}\label{eq:incl}
 -Y^++d_\alpha\in\partial f(L^+),\qquad -Y^+\in\partial g(Z^+),
\end{equation}
其中
\begin{equation}\label{eq:dalpha}
 d_\alpha=\rho\bigl((\alpha-1)a+\Delta Z\bigr)
 =\rho\bigl((\alpha-1)r+(2-\alpha)\Delta Z\bigr).
\end{equation}
\end{proposition}
\begin{proof}
第一个子问题的最优性条件是
$-Y-\rho(L^++Z-D)=-Y-\rho a\in\partial f(L^+)$。
由 $Y^+=Y+\rho(\alpha a+\Delta Z)$，可将左端写成
$-Y^++\rho((\alpha-1)a+\Delta Z)$。第二个子问题则直接给出
$-Y-\rho(\widehat L+Z^+-D)=-Y^+\in\partial g(Z^+)$。
最后用 $a=r-\Delta Z$ 化简，即得第二种表达式。
\end{proof}

当 $\alpha=1$ 时，$d_\alpha=\rho\Delta Z$，与通常写法一致。当 $\alpha\ne1$ 时，只检查 $\rho\Delta Z$ 会漏掉含 $r$ 的项。原项目的一个历史例子中，按该朴素量可在第 380 步停止，但正确缺陷尚约为容差的 $2.4$ 倍；正确的双残差条件至第 656 步才同时满足。这是停止规则的改变，不能简单解释为过松弛“加速失败”。

若近端由部分谱近似求解，第一条包含关系还带有该近端误差，$d_\alpha$ 本身便不再是完整的一阶缺陷。因此原程序把迭代中的这一量标为诊断量是必要的。罚参数改变时，还必须同步取 $u\leftarrow(\rho_{\mathrm{old}}/\rho_{\mathrm{new}})u$，以保持未缩放乘子 $Y$ 连续。有限预热后冻结参数可转入固定参数分析；无限次自适应更改需要另行验证收敛假设，不能只援引固定参数 ADMM 的结论\cite{admm}。

\subsection{完整近端固定点检查}
取独立于算法罚参数的单位近端步长。对任一候选 $(L,Z,Y)$，定义
\begin{equation}\label{eq:fixed}
 q_L=\fro{L-\prox_f(L-Y)},\qquad
 q_Z=\fro{Z-\prox_g(Z-Y)},\qquad r=L+Z-D.
\end{equation}
由 $x=\prox_f(x-Y)$ 当且仅当 $-Y\in\partial f(x)$，可知 $r=0,q_L=q_Z=0$ 恰是 KKT 条件。它们足以保证全局最优：对任意可行 $(L',Z')$，两个子梯度不等式相加给出
$f(L')+g(Z')\geq f(L)+g(Z)-\ip Y{(L'-L)+(Z'-Z)}=f(L)+g(Z)$。

反过来，在本文有限处处的范数模型中，等式约束的常规资格条件成立，最优解可配上这样的乘子。数值上，$q_L,q_Z$ 很小只说明相应固定点条件近似满足；若没有误差界常数或有界比较集合，不能仅由它们给出任意声称的目标相对误差。

原 \texttt{validate\_final\_results.m} 使用完整 SVD 重算 \eqref{eq:fixed}，并检查原始残差、分组变量一致性及元信息。其代码的判定阈值含 $1.05$ 的余量，应与记录中的实际残差比例一并报告。历史表中 8 个方法均通过。本次扩展私人资料后，另对 PCP、Relaxed、Adaptive、Partial、Combined、Stable、Gram、StableGram 八组保存端点重新核验：按原 MATLAB 灰度系数和列主序重建 200 帧输入，使用 SciPy/LAPACK 重新计算完整薄 SVD，并重算第二块近端与噪声分组一致性。八组均与历史 CSV 相符，且均通过原有 $1.05$ 倍容差的判定。观测到的最大原始可行性、低秩近端、第二块近端及噪声一致性相对容差比分别约为 $0.9980$、$0.6108$、$2.34\times10^{-11}$、$2.59\times10^{-11}$；这些最大值不一定来自同一方法。\par
原始矩阵逐项绑定 SHA256，匿名数值记录见 \texttt{evidence/endpoint\_checks.json}。本次检查使用了保存的 $L,S,E,Y$，不是重新优化、重新计时或重验每个中间迭代；完整矩阵和原画面仍留在私人归档。它也仍是浮点检查，不是区间数值认证。这 8 组端点复算与第 7 节的 222 项独立合成检查是两套分别统计的证据。

\subsection{PCP 的可行上界与对偶下界}
PCP 的对偶可写为
\begin{equation}\label{eq:dual}
 \max_Q\ \ip QD\quad\text{s.t.}\quad \op Q\leq1,\quad\maxnorm Q\leq\lambda.
\end{equation}
这来自核范数与算子范数、$\ell_1$ 与 $\ell_\infty$ 的对偶关系：在 $Q$ 满足上述约束时，任意可行 $(L,S)$ 都有
$\ip QD=\ip QL+\ip QS\leq\nuc L+\lambda\ellone S$。
给定任意乘子估计 $Y$，令
\begin{equation}
 Q=\frac{-Y}{\max\{1,\op Y,\maxnorm Y/\lambda\}},\qquad
 U=\nuc L+\lambda\ellone{D-L},\quad B=\ip QD.
\end{equation}
于是严格算术下 $B\leq v_\lambda(D)\leq U$。$U$ 特意使用可行修正 $S=D-L$。直接把尚未可行的迭代对 $(L,S)$ 的目标与 $B$ 比较，可能得到看似“违反对偶性”的数值。

原 Python 核验日志给出 $B\approx2216.3557$、$U\approx2216.3840$，而未修正迭代目标约为 $2216.3451<B$；这不构成矛盾。$(U-B)/U\approx1.3\times10^{-5}$ 是以可行上界为分母的相对间隙。若要改成以最优值为分母，应在 $B>0$ 时使用 $(U-B)/B$ 的上界。实际浮点计算未做区间包络，因此应称“数值原始—对偶夹逼”，而非机器认证的严格区间。


\subsection{近似参考点下的下降量核验}
原项目还核验了固定参数、$\alpha=1$ 的经典 ADMM 下降量。若 $(\bar L,\bar S,\bar Y)$ 是精确 KKT 点，则
\begin{equation}
 V_k=\rho^{-1}\fro{Y_k-\bar Y}^2+\rho\fro{S_k-\bar S}^2
\end{equation}
满足 $V_{k-1}-V_k\geq\rho\fro{r_k}^2+\rho^{-1}\fro{d_k}^2$，其中 $d_k=\rho(S_k-S_{k-1})$。然而数值实验得到的参考点也只是有限步迭代点，不能直接把它当作精确 KKT 点。

这个误差可以显式计入。假定当前与参考点的两个近端子问题精确求解，故
\begin{equation}
 -\bar Y+\bar d\in\partial f(\bar L),\quad-\bar Y\in\partial g(\bar S),\quad
 \bar r=\bar L+\bar S-D,
\end{equation}
当前点有对应的 $(r_k,d_k)$ 条件；并假设旧点也满足 $-Y_{k-1}\in\partial g(S_{k-1})$。则
\begin{equation}\label{eq:refbudget}
 V_{k-1}-V_k\geq\rho\fro{r_k}^2+\rho^{-1}\fro{d_k}^2-E_k,
\end{equation}
其中可取
\begin{align}
 E_k=2\bigl[&(\sqrt{\rho V_k}+\fro{d_k})\fro{\bar r}\nonumber\\
 &+\fro{\bar d}(\sqrt{V_k/\rho}+\fro{r_k}+\fro{\bar r})\bigr].
\end{align}
为验证它，记 $A=Y_k-\bar Y$、$B=S_k-\bar S$。将两个子微分的单调性相加，得到
\begin{equation*}
 \ip A{r_k}+\ip B{d_k}
 \leq\ip A{\bar r}+\ip{r_k-\bar r}{d_k}
       -\ip{L_k-\bar L}{\bar d}.
\end{equation*}
另一方面，展开两个平方差有
\begin{equation*}
 V_{k-1}-V_k=-2\ip A{r_k}-2\ip B{d_k}
             +\rho\fro{r_k}^2+\rho^{-1}\fro{d_k}^2.
\end{equation*}
相邻第二块子梯度的单调性给出 $\ip{r_k}{d_k}\leq0$。再用
$L_k-\bar L=r_k-\bar r-B$，以及
$\fro A\leq\sqrt{\rho V_k}$、$\fro B\leq\sqrt{V_k/\rho}$，对剩余内积应用 Cauchy--Schwarz 即得 \eqref{eq:refbudget}。精确参考点的情形令 $\bar r=\bar d=0$ 即可。原零初值在范数第二块上满足所需旧点条件；一般初始化应从该条件已成立的一步开始计数。

历史程序用 4000 步参考点检查前 1000 步，参考原始、对偶残差约为 $5.841\times10^{-4}$ 和 $6.242\times10^{-7}$。其最后一步预算 $E_k\approx5.556\times10^{-4}$，所记录下降余量约为 $4.393\times10^{-3}$，全部受检步的余量与预算之比最小约为 $7.9$。因此原工作确实考虑了参考点不精确的问题。这仍是有限轨道、浮点近端的数值检查；若近端有额外误差，上述预算还须加入相应项，不能借该表代替一般收敛证明或认证浮点近端的精确性。

\section{严格重复列的模型尺度与迭代尺度}
\subsection{最优值的精确缩放关系}
对正整数 $k$，定义横向重复算子
$\rep_k(A)=[A,\ldots,A]\in\R^{m\times kn}$。首先注意
\begin{equation}\label{eq:repeatnorm}
 \nuc{\rep_k(A)}=\sqrt{k}\,\nuc A,\qquad
 \ellone{\rep_k(A)}=k\ellone A.
\end{equation}
第二式是逐元素求和。第一式由
$\rep_k(A)\rep_k(A)^T=kAA^T$ 得到：每个非零奇异值乘以 $\sqrt{k}$。

\begin{proposition}[严格重复列的 PCP 最优值]\label{prop:repeat}
对任意 $D$、$\lambda>0$ 和正整数 $k$，有
\begin{equation}\label{eq:repeatvalue}
 v_\lambda(\rep_k(D))=\sqrt{k}\,v_{\sqrt{k}\lambda}(D).
\end{equation}
任何短问题权重为 $\sqrt{k}\lambda$ 的最优分解，经重复后都是长问题权重为 $\lambda$ 的最优分解。反向地，任何长问题最优分解都可以通过块置换平均得到一个块相同的最优分解，其平均块是相应短问题的最优分解。这里没有断言长问题所有最优解都必须块相同。
\end{proposition}
\begin{proof}
设长问题的任意可行分解为 $L=[L_1,\ldots,L_k]$、$S=[S_1,\ldots,S_k]$，则 $L_j+S_j=D$。同时对 $L,S$ 作相同的列块置换仍然可行，且不改变两种范数。对全部 $k!$ 个置换取平均，得到
$\rep_k(\bar L),\rep_k(\bar S)$，其中 $\bar L=k^{-1}\sum_jL_j$、$\bar S=k^{-1}\sum_jS_j$。由凸性，平均后的目标不超过原目标。再用 \eqref{eq:repeatnorm}，有
\begin{align}
 \nuc L+\lambda\ellone S
 &\geq \sqrt{k}\,\nuc{\bar L}+k\lambda\ellone{\bar S}\nonumber\\
 &=\sqrt{k}\bigl(\nuc{\bar L}+\sqrt{k}\lambda\ellone{\bar S}\bigr)
 \geq\sqrt{k}\,v_{\sqrt{k}\lambda}(D).
\end{align}
对所有长问题可行分解取下确界给出一个方向。另取短问题的最优分解 $(L_*,S_*)$，重复后可行，且目标恰为 $\sqrt{k}v_{\sqrt{k}\lambda}(D)$，给出反向不等式。
若从长问题最优解出发，上述所有不等式总的两端相等，因此平均后的目标也最优，平均块的短问题目标也必须最优。证明未使用严格凸性，故不能进一步推出未平均的每个块都相同。
\end{proof}

命题说明，增加完全重复的数据列，即使没有增加新的信息，也会改变两个正则项的相对比例。固定长问题的 $\lambda$ 相当于把短问题的权重改为 $\sqrt{k}\lambda$。若希望重复前后表示同一个短问题权重，应取 $\lambda_{\mathrm{long}}=\lambda_{\mathrm{short}}/\sqrt{k}$。

常见规则 $\lambda=1/\sqrt{\max(m,n)}$ 还需区分尺寸区间。当 $kn\leq m$ 时，重复前后的默认 $\lambda$ 都是 $1/\sqrt m$，有效短问题权重却增长为 $\sqrt{k/m}$；当 $kn\geq m$ 时，长问题默认权重为 $1/\sqrt{kn}$，折回短问题后为 $1/\sqrt n$。这个推导只针对严格重复；真实长视频中新出现的运动、光照变化和停留结构不能被该等式抹去。原报告第 32 页借重复比例作尺度解释，可以被本命题支持其理想化部分，不能据此证明一般长视频的经验调参公式。

\subsection{同一轨道需要同时变换罚参数}
\begin{proposition}[精确 ADMM 的重复协变关系]\label{prop:trajectory}
固定 $\rho>0$ 和同一松弛参数 $\alpha$，只考虑 PCP、精确近端以及固定参数迭代。长问题使用数据 $\rep_k(D)$ 和参数 $(\lambda,\rho)$；短问题使用数据 $D$ 和参数 $(\sqrt{k}\lambda,\sqrt{k}\rho)$。如果长问题初始 $(L,Z,u)$ 是短问题对应初值的逐块重复，则每一步长问题的 $(L,Z,u)$ 仍是短问题对应变量的重复。未缩放乘子的对应关系为 $Y_{\mathrm{short}}=\sqrt{k}\,Y_{\mathrm{long,block}}$。
\end{proposition}
\begin{proof}
先由重复矩阵的 SVD 或命题 \ref{prop:gram} 得到
\begin{equation}\label{eq:svtrepeat}
 \SVT_t(\rep_k(A))=\rep_k(\SVT_{t/\sqrt{k}}(A)).
\end{equation}
因此长问题的 $L$ 更新阈值 $1/\rho$ 对应短问题阈值 $1/(\sqrt{k}\rho)$。稀疏更新的软阈值在两边相同，因为
$(\sqrt{k}\lambda)/(\sqrt{k}\rho)=\lambda/\rho$；逐元素软阈值与重复算子可交换。松弛组合和 $u$ 更新都是线性组合，也与重复可交换。由初值归纳即得全部迭代的结论。最后 $Y=\rho u$，故未缩放乘子多一个 $\sqrt{k}$ 因子。
\end{proof}

这个结果并不意味着两个程序会在同一步自动停止。残差的维度尺度、绝对容差项以及参数相关的对偶容差也随问题变化；若实现的停止规则没有相同的缩放性质，同轨道的两个程序可以截断于不同步数。部分 SVD、浮点误差和自适应罚参数也超出该精确协变命题的条件。

本命题是本次整理根据原报告的重复尺度讨论进一步写明的推导，不追溯为原项目已经证明的结果，也不声称该初等等价关系具有新的文献优先权。它的实际用途是设计一致性测试：在关闭停止判断的固定步数试验中，检查参数变换后的两条轨道是否符合恒等式。

\section{密集噪声、Huber 消元与 BB 步长}
\subsection{消元后的模型与分组近端}
原项目还比较了带平方惩罚的密集噪声模型
\begin{equation}\label{eq:noisy}
 \min_{L+S+E=D}\ \nuc L+\lambda\ellone S+\frac1{2\eta}\fro E^2,
 \qquad \eta>0.
\end{equation}
稳定 PCP 文献\cite{stable}为同时考虑小密集噪声和大稀疏误差提供了经典背景；其中约束型噪声界模型与 \eqref{eq:noisy} 的平方惩罚参数并非可以不加条件地逐一互换。本文只分析明确写出的 \eqref{eq:noisy}。

固定 $Z=S+E$，逐元素消去 $S$ 得到
\begin{equation}
 h_{\eta,\lambda}(z)=\min_s\left\{\lambda|s|+\frac{(z-s)^2}{2\eta}\right\}
 =\begin{cases}z^2/(2\eta),&|z|\leq\eta\lambda,\\
 \lambda|z|-\eta\lambda^2/2,&|z|>\eta\lambda.
 \end{cases}
\end{equation}
最优 $s=\soft_{\eta\lambda}(z)$，$e=z-s=\operatorname{clip}(z,-\eta\lambda,\eta\lambda)$。于是 \eqref{eq:noisy} 等价于
$\min_{L+Z=D}\{\nuc L+H_{\eta,\lambda}(Z)\}$，其中 $H$ 是逐元素 $h$ 之和。

对 ADMM 第二块，标量近端有闭式
\begin{equation}\label{eq:huberprox}
 \prox_{h_{\eta,\lambda}/\rho}(v)
 =\begin{cases}
 \dfrac{\rho\eta}{1+\rho\eta}v,&|v|\leq\lambda(\eta+1/\rho),\\
 v-(\lambda/\rho)\sgn(v),&|v|>\lambda(\eta+1/\rho).
 \end{cases}
\end{equation}
其证明可直接求解 $h'(z)+\rho(z-v)=0$：二次区给出第一式，将它代入 $|z|\leq\eta\lambda$ 得到对应的 $v$ 区间；外侧给出第二式及其一致的区间条件。两式在边界连续。
原实现先算
\begin{equation}
 S=\soft_{\lambda(\eta+1/\rho)}(v),\qquad
 E=\frac{\rho\eta}{1+\rho\eta}(v-S),\qquad Z=S+E,
\end{equation}
与 \eqref{eq:huberprox} 完全相同。它是二块分组近端的闭式求解，并不是把三个变量未经分析地逐个循环的三块 ADMM。

\subsection{固定 Huber 分区上的两个 BB 步长}
也可以将 $Z$ 消去，对 $F(L)=\nuc L+f(L)$ 作近端梯度，其中
$f(L)=H_{\eta,\lambda}(D-L)$。其梯度逐元素为
$[\nabla f(L)]_{ij}=\operatorname{clip}((L_{ij}-D_{ij})/\eta,-\lambda,\lambda)$，Lipschitz 常数不超过 $1/\eta$。标准固定步长 $\eta$ 因而有下降保证。原项目比较的谱步长近端梯度属于已有 SpaRSA 等方法的思路\cite{sparsa}，不是新的 BB 定义。

取两点 $L_0,L_1$，向量化后令 $s=\operatorname{vec}(L_1-L_0)$、$y=\operatorname{vec}(\nabla f(L_1)-\nabla f(L_0))$。当 $s^Ty>0$ 时，定义
\begin{equation}
 t_{\mathrm{BB1}}=\frac{s^Ts}{s^Ty},\qquad
 t_{\mathrm{BB2}}=\frac{s^Ty}{y^Ty}.
\end{equation}
\begin{proposition}[Huber 割线的条件式解释]\label{prop:bb}
对 $\eta,\lambda>0$，只要 $s^Ty>0$，均有
\begin{equation}\label{eq:bbgeneral}
 \eta\leq t_{\mathrm{BB2}}\leq t_{\mathrm{BB1}}.
\end{equation}
若每个坐标从 $L_0$ 到 $L_1$ 的整段都位于同一闭二次区，或者位于同一个仿射侧区，记二次区坐标集合为 $Q$。当 $\|s_Q\|>0$ 时进一步有
\begin{equation}\label{eq:bbexact}
 t_{\mathrm{BB2}}=\eta,\qquad
 t_{\mathrm{BB1}}=\eta\,\frac{\|s\|^2}{\|s_Q\|^2}\geq\eta.
\end{equation}
边界点按整段所在的区域归类；位于边界且位移为零的坐标不影响任一等式。
\end{proposition}
\begin{proof}
梯度的每个标量分量单调且 $1/\eta$-Lipschitz，故可写 $y_i=d_i s_i$，其中 $0\leq d_i\leq1/\eta$；$s_i=0$ 时取 $d_i=0$。于是
$\|y\|^2=\sum_i d_i^2s_i^2\leq\eta^{-1}\sum_i d_is_i^2=\eta^{-1}s^Ty$，给出第一不等式。Cauchy--Schwarz 给出 $(s^Ty)^2\leq\|s\|^2\|y\|^2$，给出第二不等式。固定分区条件下，二次区 $y_i=s_i/\eta$，仿射区 $y_i=0$，故
$s^Ty=\|s_Q\|^2/\eta$、$y^Ty=\|s_Q\|^2/\eta^2$，代入即可。
\end{proof}

这解释了原实验中 BB2 常停留在接近 $\eta$ 的尺度，而 BB1 可以取大得多的步长；它没有证明 BB1 总是更快。若割线跨过 Huber 的转折点，只有 \eqref{eq:bbgeneral} 一般成立，不能再写 $t_{\mathrm{BB2}}=\eta$。近端映射、非单调线搜索、步长上限、起点和目标精度也都会改变实际调用次数。

原报告还用支撑稳定解释谱罚参数为何在一个 PCP 例子中不触发。可以严格保留的部分是：若 $-Y_j\in\lambda\partial\|S_j\|_1$，且某坐标的两次 $S$ 都非零同号，则该坐标 $\Delta Y=0$；若两次都为零，则 $\Delta S=0$。因此 $\ip{\Delta S}{-\Delta Y}$ 的贡献只来自支撑或符号发生变化的坐标。这并不证明整个问题没有可用的谱信息，更不能否定处理非光滑问题的谱 ADMM 文献\cite{aadmm}。观测到相关性阈值不通过，只能作为该次轨道的机制证据。

\section{历史实验逐组能支持什么}
\subsection{保持同一模型的实现比较}
以下数字来自原项目记录，不是本次在同一机器重新计时。真实视频矩阵尺寸为 $25344\times200$，没有像素级真实分解。默认 PCP 在双残差容差下运行 608 步；固定参数 Gram 实现也是 608 步，最终背景相对差约为 $1.21\times10^{-14}$。原单次用时分别约为 $128.2$ 秒和 $85.0$ 秒，SVT 部分分别约为 $81.2$ 秒和 $31.0$ 秒。它支持该输入上输出一致与一次运行的时间差，不能直接外推到其他尺寸或平台。

12 次核函数比较的时间中位数之比约为 $2.16$。完整求解器又作了三组配对复核，原始数据为：
\begin{center}
\begin{tabular}{rrrr}
\toprule
配对 & 完整 SVD／秒 & Gram／秒 & 完整 SVD 时间／Gram 时间\\
\midrule
1 &177.2119&99.9946&1.7722\\
2 &96.1714&609.1713&0.1579\\
3 &100.4218&69.3531&1.4480\\
\bottomrule
\end{tabular}
\end{center}
全部配对的残差、输出一致性检查均通过。第二组 Gram 轨道有一次约 518 秒的异常单步耗时，现有材料没有独立解释其原因。保留原值时，配对比值中位数为 $1.4480$，范围 $0.1579$--$1.7722$；分别取两种方法时间中位数后再相除则为 $1.0043$。这两种统计量回答不同问题，不能混用。原报告中剔除该单步后的比较只能视为事后敏感性分析，不能替代原始计时。综上，当前证据支持核函数有加速潜力，不支持“完整求解器稳定提速两倍”。

部分 SVD 的历史版本同样运行 608 步且无回退，但总用时约为 181 秒，比完整 SVD 慢。这个负面结果与原问题列数仅 200、迭代末保留秩接近 91 的结构相容；在这样的维度和秩比例下，反复构造部分谱未必节省工作。Gram 单元表的 11 项检查中，最大相对输出差约为 $5.09\times10^{-15}$，4 项触发回退；这些覆盖了若干秩亏、近阈值和极端量级输入，不构成对所有浮点输入的保证。

\subsection{松弛与罚参数必须在同一精度口径下比较}
原 PCP 与 $\alpha=1.6$ 的松弛版本，在完整双残差停止下分别为 608 和 656 步。若仅比较达到同一原始残差目标的步数，原始残差阈值按既定容差的倍数依次取 $1,0.8,0.6,0.5,0.415$，两者分别为
\begin{equation*}
 (608,699,838,937,1049),\qquad(380,438,524,586,656).
\end{equation*}
后者约为前者的 $1/1.6$。这说明可行性沿该轨道进展更快，但不代表在相同完整 KKT 精度下也有相同倍数的改进。命题 \ref{prop:relax} 给出了两种说法并存的原因。

固定罚参数 $\rho=0.05,0.1,0.2,0.4,1$ 的记录分别为 $1211,608,310,211,413$ 步，表明默认 $0.1$ 不是这个样例中最好的已试参数。原有预热残差平衡规则未触发任何修改，故与基线相同；将平衡阈值调低后有 411 步和 357 步的记录。残差平衡具有尺度敏感性，不能把某一阈值在一个输入上表现良好当成通用结论\cite{wohlberg}。

谱罚参数试验在该 PCP 样例中进行了 303 次相关性检查，均未通过 $0.2$ 的接受阈值，因此也是 608 步。另一带噪模型中谱规则会更新参数，但对应步数 71、93、117 并没有超过该组已试固定参数的最佳表现；这个结果也应保留。相关性筛选通过、参数确实变化与整体收益是三个需要分别报告的量。

\subsection{改变模型后的误差与计算量}
带密集噪声模型在该真实输入上的一个版本只需 26 步，Gram 版本的总记录约为 3.4 秒。但它求解的是 \eqref{eq:noisy}，不是 \eqref{eq:pcp}，故“608 步降到 26 步”包含模型改变，不能作为同一目标的纯算法加速倍数。真实数据没有真实分量，前景面积、时间中位数代理和视觉观察只能描述输出行为，不能直接给出恢复准确率。

有真实分量的合成组包含随机低秩与稀疏扰动、移动矩形、静止矩形及不同噪声强度，共 60 行主结果。静止亮块可被吸入低秩项，PCP 的该组前景误差约为 $0.896$、F1 约为 $0.447$；这属于模型识别失败，与先前矩形证书专题重叠。它不应被删去以强化方法效果。

在同一生成族、三个留出种子和噪声标准差 $0.03$ 的比较中，PCP 的平均前景相对误差约为 $0.4325$，带噪模型约为 $0.2343$；背景误差约由 $0.0173$ 到 $0.0153$，平均步数约由 249 到 31。未知噪声试验另有两种子、两噪声强度、两模型共 8 行记录，估计标准差与真值比约为 $1.02$--$1.03$。它们是有用的同族留出检查，仍不是外部数据集验证。

误差指标也会相互冲突。例如某带噪参数下前景误差改善而背景误差略增；二值化 F1 高并不保证前景幅值正确。18 行阈值扫描与 30 行参数敏感性记录应和主表一起阅读。其中参数敏感性表有 2 行在 1500 步上限内未收敛，不能按已收敛结果解释其误差。

\subsection{长序列、秩与模型失配}
1000 帧合成长序列的默认记录为 2134 步，保留秩 73，背景和前景相对误差约为 $0.102$ 和 $0.597$。本次另对已保存的 1000 个不同合成帧及其真实分量独立复算，背景、前景相对误差为 $0.1021635$、$0.5967037$，F1 为 $0.8060540$，完整 SVD 重算的低秩近端残差为 $4.75417\times10^{-6}$，与历史结果相符。这项核验比对保存端点与真实分量，没有重新运行优化、改变参数或重新计时；匿名记录见 \texttt{evidence/synthetic\_1000\_endpoint.json}。\par
即使完整近端检查通过，这也不是正确分解的证据。原表给出的某个可行分解目标约为 $1433.5$，真实分解目标约为 $1471.5$；只要这两个目标的计算可靠，前者更小已经足以排除真实分解的最优性，无须先证明数值解达到最优。这里的目标差约为 $2.6\%$，远大于所示末位舍入，但当前仍是历史浮点证据，不是区间认证。

将权重依次乘以 $1,0.8409,0.7071,0.5,0.3448$ 后，历史步数为 $2134,902,50,38,33$，保留秩为 $73,35,3,3,3$。它说明这个生成实例对权重很敏感。命题 \ref{prop:repeat} 可以解释严格重复的部分尺度机制，却不能将这里所有真实不同帧化为严格重复。另有 $17\times13\times3=663$ 行停留时间—权重—种子扫描，采用更大的 5000 步上限及不同罚参数；不能把这组“全部收敛”倒推为前述 1500 步敏感性组也已收敛。22 行理想矩形阈值核验与已有矩形专题存在实质重叠，应作为其原始证据归档。

最后，代数秩、算法阈值保留秩和能量秩是不同量。若 $\sum_{i\leq r_\varepsilon}\sigma_i^2\geq(1-\varepsilon)\sum_i\sigma_i^2$，则 $r_\varepsilon$ 只度量主要能量维数。一个保留秩约 91 的背景仍可有几乎全部能量集中在第一奇异方向。时间中心化 $L(I-\mathbf1\mathbf1^T/n)$ 又移除了整体静态背景，得到的谱回答的是变化部分的维数。原项目提供了原始谱与中心化谱，二者应并列解释，不能择一作为“真正低秩”的证明。

\subsection{课程方法比较中的计算口径}
原 Python 组比较了次梯度、固定步长近端梯度、BB1/BB2 近端梯度和 ADMM。在带噪目标的一组零初值记录中，以近似高精度参考目标衡量相对目标差 $10^{-6}$ 时，固定步长需约 201 次 SVT，BB1 为 39 次（36 次接受更新和 3 次回溯），BB2 为 181 次；改变初值后 BB1 又降到 23 次 SVT。原报告按每十步检查的固定步长表出现 210，应与逐步日志的 201 区分。若用“迭代步数”遗漏回溯中的额外 SVT，就会改变比较。

参考值来自较严格容差的数值 ADMM，而非该带噪模型的形式认证最优值。因此这些数字是相对于该参考解的经验精度，不是严格复杂度界。命题 \ref{prop:bb} 给出 BB1/BB2 差异的一种局部解释；线搜索窗口、步长上限和起点的敏感性仍需记录。它不能由一个课程算例推出某算法普遍优于另一个算法。

\section{本次独立合成核验与可重复范围}
随附 \texttt{verify\_equivalences.py} 为重新编写的 NumPy 程序，不导入原组代码，不读取视频。固定种子为 20261125，当前环境为 Python 3.12.4、NumPy 1.26.4；程序仅在环境未预先设定时请求单线程执行。输出 \texttt{evidence/synthetic\_checks.json} 保存每项误差、容差及通过状态，共 222 项断言全部通过，内容如下。
\begin{enumerate}
\item 高、宽和方阵，秩亏、零阈值、近阈值、零矩阵、极端量级上的完整 SVD 与独立保护型 Gram 重构比较；最大相对差约为 $1.14\times10^{-15}$。新保护规则只用于合成核验，不冒充原 MATLAB 保护公式，也仍是启发式。
\item 例 \ref{ex:partial} 的零三元组残差与恰为 $4$ 的近端误差；Huber 分组表达式与闭式近端、标量一阶条件一致性。
\item $\alpha=0.5,1,1.6,1.95$ 下 PCP 与 Huber 两种第二块的单步子梯度恒等式，以及 \eqref{eq:dalpha} 两种表达式。
\item 固定分区的 BB 两等式和 100 组跨分区的界 \eqref{eq:bbgeneral}。
\item $k=2,3,7$ 的重复范数、块平均不增目标、固定 40 步的 ADMM 轨道关系，以及分别求解短、长问题所得最优值的数值夹逼。40 步轨道的最大相对差约为 $5.48\times10^{-14}$。
\end{enumerate}

独立求解试验使用 $20\times8$ 的合成矩阵及其重复。它没有用命题 \ref{prop:trajectory} 的罚参数变换强制两边同轨道，而是两边都取 $\rho=0.8$，分别达到原始相对残差 $10^{-9}$ 及相对可行间隙 $10^{-8}$ 的停止条件。所得最优值对比如下。
\begin{center}
\begin{tabular}{rrrr}
\toprule
$k$ & 短问题步数 & 长问题步数 & 缩放后可行上界相对差\\
\midrule
2 &1580&660&$7.78\times10^{-13}$\\
3 &199440&116240&$1.57\times10^{-12}$\\
7 &920&360&$2.82\times10^{-10}$\\
\bottomrule
\end{tabular}
\end{center}
$k=3$ 的试验在首次设置的 20000 步上限内没有达到精度要求：原始残差约为 $4.51\times10^{-6}$、可行间隙约为 $9.55\times10^{-6}$。最终保留同一矩阵、罚参数和停止标准，将上限提高到 200000 后完成。这里特意保留慢收敛事实；本次试验验证的是等价关系和数值检查，不是为了展示求解速度。

这些有限样例不会代替证明，也不会自动覆盖原 25344 行真实矩阵的端点；后者已用私人包中的完整矩阵另作第 3.2 节所述复算。运行程序只需 NumPy，所有生成数据驻留内存，结果写入同目录下的 \texttt{evidence}。复验时若某环境在最大步数内不达停止条件，程序会报错，不会将末步静默标为通过。浮点误差、BLAS 差异可能改变最后步数；数学恒等式则由正文证明承担。

\section{归属、覆盖范围与可形成的独立报告}
原始材料属于三人课程小组项目；最终个人报告自述承担建模、实现和实验，课堂讲演又有明确的小组分工。本稿只据材料陈述来源，不作个人贡献或著作权的独立鉴定，不将全部原始内容改署为某一人的工作。公开随附的合成核验程序重新生成测试矩阵；真实保存端点仅使用已归档矩阵在私人资料中复算，本文报告匿名数值摘要，不随附原始视频、隐私数据、原组代码或身份信息。

从学术内容看，可以分成三层。PCP、SVT、ADMM、Huber 消元、谱罚参数、非单调谱近端梯度及原始—对偶界都属于已有方法或基本推论。原项目可明确保留的工作是面向具体矩阵形状的实现选择、完整端点复核、停止规则纠正、参数实验及模型失败诊断。本文把严格重复下的最优值关系、相应轨道变换及 Huber BB 条件式解释写成完整命题，目的是使这些讨论可以独立检查；在没有专门系统先行文献检索前，不把它们命名为新定理或新方法。

相较于已经整理的静止矩形证书稿，本稿的独立价值在于“等价计算与核验”：同一近端的不同实现必须先比较输出；加速叙述必须固定模型和精度；可行性与一阶条件必须分别检查；求解成功也不能代替模型正确。现有材料足够支撑这篇带证明的数值方法阅读报告，尚不足以支撑广泛性能优势、真实场景恢复准确率或经认证的浮点稳定性结论。

\clearpage
\appendix
\section{实验与程序阅读映射}
下表记录实验组及其主要证据；原件留在私有归档。原课程目录与交付 ZIP 各有 326 个文件，逐相对路径和 SHA256 一致。本稿另存 102 项匿名来源哈希，统计 48 张历史 CSV 表。文件名均为不含身份的相对名；完整列名、行数和状态计数见匿名证据文件。
\small
\setlength{\LTpre}{0pt}\setlength{\LTpost}{0pt}
\begin{longtable}{>{\raggedright\arraybackslash}p{.23\textwidth}>{\raggedright\arraybackslash}p{.35\textwidth}>{\raggedright\arraybackslash}p{.34\textwidth}}
\toprule
组别 & 原程序／历史表 & 本稿保留的结论或边界\\
\midrule\endhead
核心求解与近端 & \texttt{rpca\_admm.m}, \texttt{svt\_gram.m}, \texttt{svt\_adaptive.m} & 二块分组、正确松弛缺陷；Gram 保护和部分谱检查均不等于严格遗漏谱认证。\\
最终端点复核 & \texttt{validate\_final\_results.m}, \texttt{final\_validation.csv}（8 行） & 8 组端点已从独立私人包复算通过；不等于重跑优化。\\
基本合成与真实组 & \texttt{run\_experiments.m}, \texttt{synthetic.csv}（60）、\texttt{real.csv}（8） & 有真值与无真值指标分开；模型改变与实现改变分开。\\
参数敏感性与留出 & \texttt{run\_sensitivity.m}, \texttt{run\_holdout\_noise.m}；30 与 6 行 & 2 项到上限未收敛；留出仅在同一生成族内。\\
未知噪声及阈值 & \texttt{run\_review\_unknown\_noise.m}；8 与 18 行 & 估计代价和二值化阈值不能省略。\\
长序列与模型诊断 & \texttt{run\_scale\_video.m}, \texttt{analyze\_scale\_result.m}；\texttt{scale\_video.csv}, \texttt{scale\_analysis.csv} & 可行目标低于真实目标足以排除真实分解最优性；并非迭代失败。\\
权重与停留时间 & \texttt{run\_lambda\_sweep.m}, \texttt{rpca\_numpy\_lambda.py}；5、663、22 行 & 理想矩形已有专题覆盖；真实长序列不能直接套严格重复等式。\\
真实权重扫描 & \texttt{lambda\_real.csv}（9），\texttt{lambda\_real\_stable.csv}（7） & 无真实分解；不公开具体区域坐标或画面。\\
Gram 单元与内核 & \texttt{test\_gram\_svt.m}, \texttt{benchmark\_svt\_kernel.m}；11 与 12 行 & 记录回退、误差与成本；小矩阵不保证更快。\\
边界与配对计时 & \texttt{test\_review\_edges.m}, \texttt{run\_review\_round1.m}；15 项、3 配对 & 原始计时异常完整保留；输出一致与计时离群分别判断。\\
谱与能量秩 & \texttt{pcp\_background\_spectra.csv}（200）及两个汇总表 & 中心化谱与原始谱衡量不同内容。\\
NumPy 公式复核 & \texttt{rpca\_numpy\_check.py}, \texttt{rpca\_numpy\_extra.py} & 另一路实现；其 Gram 近端没有 MATLAB 全套防护。\\
课程方法与 BB & \texttt{rpca\_numpy\_course.py}, \texttt{rpca\_numpy\_bbpg.py} & 计入回溯 SVT，标注近似参考目标、起点和步长上限。\\
罚参数及 Lyapunov & \texttt{rpca\_numpy\_rho.py}, \texttt{rpca\_numpy\_lyapunov.py} & 17 行罚参数表；Lyapunov 数值轨迹不是一般收敛证明。\\
绘图、载入与输出 & \texttt{load\_video\_frames.m}, \texttt{recovery\_metrics.m} 及绘图脚本 & 服务于输入与指标；本稿不复制原画面及视频导出物。\\
\bottomrule
\end{longtable}
\normalsize

\section{命题与证据的责任划分}
\begin{center}
\small
\begin{tabularx}{\textwidth}{>{\raggedright\arraybackslash}p{.24\textwidth}>{\raggedright\arraybackslash}X >{\raggedright\arraybackslash}p{.23\textwidth}}
\toprule
陈述 & 证明或证据来源 & 不包含的结论\\
\midrule
Gram 公式 & 命题 \ref{prop:gram}，经典 SVT 的代数重写 & 全输入浮点认证、新复杂度阶\\
松弛 KKT 缺陷 & 命题 \ref{prop:relax}，两个精确近端条件 & 近似子问题的误差自动消失\\
重复最优值 & 命题 \ref{prop:repeat}，置换平均与范数缩放 & 每个最优解块相同、一般长视频等价\\
重复轨道 & 命题 \ref{prop:trajectory}，固定参数归纳 & 默认停止时刻相同、自适应情形\\
Huber BB 条件 & 命题 \ref{prop:bb}，标量割线界 & 任意割线 BB2 恒等、全局加速率\\
历史性能 & 原 CSV、日志及代码相互核对 & 从初值重新优化、普遍速度优势\\
新增 222 项核验 & 独立合成程序与逐项 JSON & 矩阵证明的形式化、区间认证\\
\bottomrule
\end{tabularx}
\end{center}

\begin{thebibliography}{99}
\bibitem{pcp} E. J. Candès, X. Li, Y. Ma, J. Wright. Robust Principal Component Analysis? \emph{Journal of the ACM}, 58(3), Article 11, 2011. 原始预印本：\url{https://arxiv.org/abs/0912.3599}.
\bibitem{svt} J.-F. Cai, E. J. Candès, Z. Shen. A Singular Value Thresholding Algorithm for Matrix Completion. \emph{SIAM Journal on Optimization}, 20(4):1956--1982, 2010. \url{https://arxiv.org/abs/0810.3286}.
\bibitem{admm} S. Boyd, N. Parikh, E. Chu, B. Peleato, J. Eckstein. Distributed Optimization and Statistical Learning via the Alternating Direction Method of Multipliers. \emph{Foundations and Trends in Machine Learning}, 3(1):1--122, 2011. \url{https://web.stanford.edu/~boyd/papers/admm_distr_stats.html}.
\bibitem{stable} Z. Zhou, X. Li, J. Wright, E. Candès, Y. Ma. Stable Principal Component Pursuit. \emph{ISIT}, 2010. 原始预印本：\url{https://arxiv.org/abs/1001.2363}.
\bibitem{sparsa} S. J. Wright, R. D. Nowak, M. A. T. Figueiredo. Sparse Reconstruction by Separable Approximation. \emph{IEEE Transactions on Signal Processing}, 57(7):2479--2493, 2009. 作者公开版：\url{https://www.lx.it.pt/~mtf/SpaRSA/IEEE_TSP_2009_Wright_Nowak_Figueiredo.pdf}.
\bibitem{aadmm} Z. Xu, M. Figueiredo, T. Goldstein. Adaptive ADMM with Spectral Penalty Parameter Selection. \emph{AISTATS}, PMLR 54:718--727, 2017. \url{https://proceedings.mlr.press/v54/xu17a.html}.
\bibitem{wohlberg} B. Wohlberg. ADMM Penalty Parameter Selection by Residual Balancing. 2017. \url{https://arxiv.org/abs/1704.06209}.
\end{thebibliography}
\end{document}
